\documentclass[11pt,a4paper]{article}
\usepackage[utf8]{inputenc}
\usepackage[T1]{fontenc}
\usepackage{geometry}
\geometry{margin=1in}
\usepackage{hyperref}
\usepackage{enumitem}
\usepackage{booktabs}
\usepackage{longtable}
\usepackage{xcolor}
\usepackage{titlesec}
\usepackage{amsmath,amssymb}

\titleformat{\section}{\Large\bfseries\color{blue!70!black}}{}{0em}{}
\titleformat{\subsection}{\large\bfseries\color{blue!50!black}}{}{0em}{}
\titleformat{\subsubsection}{\normalsize\bfseries\color{blue!30!black}}{}{0em}{}

\title{\textbf{Replication Plan}\\[0.5em]
\large Chiral Spin Order in Kondo--Heisenberg Systems\\[0.3em]
\normalsize OSTI ID: 1412756}
\author{Auto-generated Replication Blueprint}
\date{\today}

\begin{document}
\maketitle
\tableofcontents
\newpage

%---------------------------------------------------------------------
\section{Paper Overview}
%---------------------------------------------------------------------
This paper investigates chiral spin order in Kondo--Heisenberg systems on triangular and kagome lattice geometries. The authors study the competition between RKKY-type indirect exchange and Kondo screening in geometrically frustrated lattices, focusing on the emergence of scalar spin chirality and its interplay with charge transport (anomalous Hall effect). The work employs mean-field decoupling of the Kondo lattice Hamiltonian, self-consistent numerical solutions, and variational approaches to map out phase diagrams as functions of Kondo coupling $J_K$, Heisenberg exchange $J_H$, and electron filling.

%---------------------------------------------------------------------
\section{Detailed Workflow}
%---------------------------------------------------------------------

\subsection{Phase 1: Hamiltonian Construction}
\begin{enumerate}[label=\textbf{1.\arabic*}]
  \item \textbf{Define the Kondo--Heisenberg Hamiltonian.}
    \begin{itemize}
      \item Tight-binding hopping term $H_t = -t \sum_{\langle i,j\rangle} c^\dagger_{i\sigma} c_{j\sigma}$ on the triangular/kagome lattice.
      \item Heisenberg exchange $H_H = J_H \sum_{\langle i,j\rangle} \mathbf{S}_i \cdot \mathbf{S}_j$ between localized spins.
      \item Kondo coupling $H_K = J_K \sum_i \mathbf{s}_i \cdot \mathbf{S}_i$ between itinerant electron spin and localized spin.
    \end{itemize}
  \item \textbf{Construct lattice geometry.}
    \begin{itemize}
      \item Build adjacency matrices for 2D triangular and kagome lattices with periodic boundary conditions.
      \item Support system sizes $N = L \times L \times n_{\text{sub}}$ with $n_{\text{sub}} = 1$ (triangular) or $3$ (kagome).
    \end{itemize}
\end{enumerate}

\subsection{Phase 2: Mean-Field Decoupling}
\begin{enumerate}[label=\textbf{2.\arabic*}]
  \item \textbf{Implement mean-field decomposition.}
    \begin{itemize}
      \item Decouple the Kondo interaction via mean-field parameters: hybridization amplitude $V_i = \langle c^\dagger_{i\sigma} f_{i\sigma} \rangle$ and bond order parameters.
      \item Introduce chiral order parameter $\chi_{ijk} = \mathbf{S}_i \cdot (\mathbf{S}_j \times \mathbf{S}_k)$ for triangular plaquettes.
    \end{itemize}
  \item \textbf{Derive the mean-field Hamiltonian.}
    \begin{itemize}
      \item Express the effective quadratic Hamiltonian in momentum space via Fourier transform.
      \item Diagonalize $H_{\text{MF}}(\mathbf{k})$ for each $\mathbf{k}$-point in the Brillouin zone.
    \end{itemize}
\end{enumerate}

\subsection{Phase 3: Self-Consistent Solution}
\begin{enumerate}[label=\textbf{3.\arabic*}]
  \item \textbf{Implement self-consistency loop.}
    \begin{itemize}
      \item Initialize mean-field parameters (hybridization, magnetization, chirality).
      \item Diagonalize the Hamiltonian at each $\mathbf{k}$-point.
      \item Compute ground-state expectation values of all order parameters.
      \item Update parameters and iterate until convergence ($\Delta < 10^{-6}$).
    \end{itemize}
  \item \textbf{Scan parameter space.}
    \begin{itemize}
      \item Vary $J_K/t$ and $J_H/t$ over specified ranges.
      \item Vary electron filling $n_c$.
      \item Record converged order parameters at each point.
    \end{itemize}
\end{enumerate}

\subsection{Phase 4: Phase Diagram Construction}
\begin{enumerate}[label=\textbf{4.\arabic*}]
  \item \textbf{Identify phases.}
    \begin{itemize}
      \item Kondo singlet phase: $V \neq 0$, $\chi = 0$.
      \item Chiral spin liquid/order: $\chi \neq 0$.
      \item Magnetically ordered phases: finite sublattice magnetization.
    \end{itemize}
  \item \textbf{Map phase boundaries} as functions of $J_K/t$, $J_H/t$, and filling.
  \item \textbf{Reproduce phase diagrams} (Figures in the paper).
\end{enumerate}

\subsection{Phase 5: Transport and Response Functions}
\begin{enumerate}[label=\textbf{5.\arabic*}]
  \item \textbf{Compute anomalous Hall conductivity} $\sigma_{xy}$ via the Kubo formula or Berry curvature integration.
  \item \textbf{Calculate density of states} and spectral functions.
  \item \textbf{Reproduce transport plots} from the paper.
\end{enumerate}

\subsection{Phase 6: Validation}
\begin{enumerate}[label=\textbf{6.\arabic*}]
  \item Compare phase boundaries with paper's figures.
  \item Compare numerical values of order parameters at specific parameter points.
  \item Verify limiting cases (e.g., $J_K = 0$ reduces to pure Heisenberg model).
\end{enumerate}

%---------------------------------------------------------------------
\section{Datasets}
%---------------------------------------------------------------------
This is a purely theoretical/computational study. No external experimental datasets are required.

\begin{longtable}{p{4.5cm} p{3.5cm} p{6cm}}
\toprule
\textbf{Dataset / Source} & \textbf{Type} & \textbf{Location / Notes} \\
\midrule
\endfirsthead
\toprule
\textbf{Dataset / Source} & \textbf{Type} & \textbf{Location / Notes} \\
\midrule
\endhead
Lattice geometries (triangular, kagome) & Synthetic & Generated programmatically from lattice vectors \\
\midrule
Model parameters ($J_K$, $J_H$, $t$, filling) & Synthetic & Specified in paper text and figure captions \\
\midrule
Paper figures (for validation) & Reference & Extracted from the PDF for quantitative comparison \\
\bottomrule
\end{longtable}

%---------------------------------------------------------------------
\section{Models, Tools, and Software}
%---------------------------------------------------------------------

\subsection{Programming Languages}
\begin{itemize}
  \item \textbf{Python 3.8+} or \textbf{Julia 1.6+} --- for implementing the mean-field solver.
  \item \textbf{Fortran/C++} --- optional, for performance-critical diagonalization loops.
\end{itemize}

\subsection{Libraries and Frameworks}
\begin{longtable}{p{4.5cm} p{2.5cm} p{6.5cm}}
\toprule
\textbf{Tool / Library} & \textbf{Version} & \textbf{Purpose / URL} \\
\midrule
\endfirsthead
\toprule
\textbf{Tool / Library} & \textbf{Version} & \textbf{Purpose / URL} \\
\midrule
\endhead
NumPy & $\geq 1.21$ & Array operations, linear algebra \\
\midrule
SciPy & $\geq 1.7$ & Sparse eigensolvers, optimization for self-consistency \\
\midrule
Matplotlib & $\geq 3.4$ & Phase diagram visualization, contour plots \\
\midrule
LAPACK/BLAS & system & Backend for matrix diagonalization \\
\midrule
PythTB & latest & Tight-binding model construction and Berry curvature \newline \url{https://www.physics.rutgers.edu/pythtb/} \\
\midrule
kwant & $\geq 1.4$ & Optional: quantum transport calculations \newline \url{https://kwant-project.org/} \\
\bottomrule
\end{longtable}

\subsection{Hardware Requirements}
\begin{itemize}
  \item Standard workstation; the mean-field calculation is per-$\mathbf{k}$-point and parallelizable.
  \item Dense $\mathbf{k}$-mesh (e.g., $100 \times 100$) requires moderate memory ($\sim$4--8\,GB).
\end{itemize}

\end{document}
